\chapter{Evolutionary Game Theory: Invasion, Dynamics and Stability}
\label{sec:p17-rat93-evolutionary-games}

Evolutionary game theory asks how strategy frequencies change when
success depends on the strategies already present. Its conclusions concern
a specified population model and adjustment rule; they need not describe
the reasoning of a utility-maximizing individual. Alexander's chapter
distinguishes resistance to invasion, evolutionary dynamics and local
stability, and shows why none alone selects a socially desirable or likely
outcome. The present reconstruction retains its complete finite treatment,
examples and qualifications, with the formal derivations in the annex.
\cite[chap.~9.3, pp.~553--564]{Alexander2021EvolutionaryGames}
\index{evolutionary game theory}\index{evolutionarily stable strategy}

The historical account begins with frequency-dependent fitness in population
biology: the fitness advantage of a behavior can change when it becomes
common. Alexander attributes the modern ESS formulation to Maynard Smith
and Price's 1973 animal-conflict paper, while noting Maynard Smith's 1972
presentation and acknowledgment of Price's unpublished work. Fisher's
sex-ratio argument is an earlier strategic motivation. These are attributed
historical connections, not claims that organisms deliberate about payoffs.
Later economic applications sought both an equilibrium-selection principle
and a theory of adjustment absent from a static equilibrium definition.
Multiple stability concepts and revision rules prevent a unique general
answer. A strategy specifying actions at every future decision node is
still a static object; it is not itself an evolutionary law of motion.
\cite[Sec.~1 and notes~1--2, pp.~553--554, 564]{Alexander2021EvolutionaryGames}

\section{Strategies, population states and resistance to invasion}
\label{sec:p17-rat93-static}

\begin{hbdefinition}{def:p17-rat93-domain}{The source's strategic and population domains}
For a finite game with players \(i=1,\ldots,N\), Alexander writes
\(S_i\) for the pure-strategy set, \(\Delta_i\) for its probability
simplex, \(\sigma_i\in\Delta_i\) for an individual mixture,
\(\vec\sigma=(\sigma_1,\ldots,\sigma_N)\) for a profile and
\(\pi_i(\vec\sigma)\) for player \(i\)'s expected payoff.
These are the objects of \cref{def:p17-finite-strategic-game}, with
\(S_i,\sigma_i,\pi_i\) taking the roles of \(A_i,x_i,U_i\).
Nash means
\[
 \pi_i(\sigma_i^*,\sigma_{-i}^*)\geq
 \pi_i(\sigma_i,\sigma_{-i}^*)\qquad(i,\ \sigma_i\in\Delta_i).
\]
The finite mixed-existence theorem is supplied by
\cref{thm:p17-nash-1950,thm:p17-nash-1951}; existence does not imply
convergence of an adjustment process.

For the single-population matrix model let \(S=\{S_1,\ldots,S_n\}\),
\(\Delta_n=\{p\in\mathbb R^n:p_i\geq0,\sum_i p_i=1\}\), and
\(A_{ij}=\pi(S_i\mid S_j)\). Role symmetry means that the payoff
pair at \((S_i,S_j)\) is \((A_{ij},A_{ji})\); it does not require
\(A_{ij}=A_{ji}\). Extend bilinearly:
\[
       \pi(\sigma\mid\mu)=\sigma^{\mathsf T}A\mu.
\]
Here \(\sigma,\mu\in\Delta_n\) can describe an individual's random
choice. In contrast, \(\vec p(t)\in\Delta_n\) describes the
frequencies of pure-strategy users at time \(t\), with no individual
randomizing in this population interpretation. Random independent
matching gives \(W_i(\vec p)=(A\vec p)_i\). The same vector can
therefore enter the same payoff algebra under two different interpretations;
the identification does not make a population share a belief or a personal
lottery. General population payoffs \(W:\Delta_n\to\mathbb R^n\)
can be nonlinear, but the ESS equivalences below use the finite linear
matrix domain. Payoff-as-fitness, random matching, the strategy repertoire
and the absence of spatial or network structure are substantive assumptions.
\end{hbdefinition}

\begin{hbdefinition}{def:p17-rat93-ess}{ESS and the two invasion comparisons}
A strategy \(\sigma\in\Delta_n\) is an evolutionarily stable
strategy (ESS) when, for every \(\mu\ne\sigma\), either
\[
 \pi(\sigma\mid\sigma)>\pi(\mu\mid\sigma),
\]
or
\[
 \pi(\sigma\mid\sigma)=\pi(\mu\mid\sigma)
 \quad\hbox{and}\quad
 \pi(\sigma\mid\mu)>\pi(\mu\mid\mu).
\]
The second comparison is required only for alternatives tied against the
incumbent. It is not a requirement that \(\sigma\) be a strict best
response to every possible opponent. A mutant population fraction
\(\varepsilon\) gives opponent distribution
\(\varepsilon\mu+(1-\varepsilon)\sigma\). Resistance to this mutant
means that some \(0<\varepsilon_\mu<1\) makes the incumbent payoff
strictly greater than the mutant payoff for every
\(0<\varepsilon<\varepsilon_\mu\). The threshold in this formulation
is allowed to depend on \(\mu\).

The strategy \(\sigma\) is locally superior when some neighborhood
of \(\sigma\) satisfies
\[
        \pi(\sigma\mid\mu)>\pi(\mu\mid\mu)
        \quad(\mu\ne\sigma).
\]
An evolutionarily stable \emph{state} is the corresponding locally
superior distribution interpreted as population frequencies. The strict
sign is essential: the zero matrix satisfies every analogous weak
inequality while having no ESS when \(n\geq2\).
\end{hbdefinition}

\begin{hbtheorem}{thm:p17-rat93-ess-equivalence}{proved=annex:p17.rat93-ess-equivalence}{Three equivalent ESS formulations}
In the stated finite bilinear domain, the two-clause ESS condition,
resistance to every sufficiently rare fixed mutant, and local superiority
are equivalent. Every ESS gives a symmetric Nash profile
\((\sigma,\sigma)\); every strict symmetric Nash strategy is an ESS.
Neither converse holds in general.
\end{hbtheorem}

\begin{hbdefinition}{def:p17-rat93-weak-dominance}{Weak dominance in the full mixed domain}
The strategy \(\mu\) weakly dominates \(\sigma\) when
\(\pi(\mu\mid q)\geq\pi(\sigma\mid q)\) for every
\(q\in\Delta_n\), with strict inequality for at least one \(q\).
Strict domination requires strict inequality for every \(q\).
The support of a distribution consists of the indices having positive
probability. These comparisons use one player's payoff scale and do not
assert interpersonal utility comparability.
\end{hbdefinition}

\begin{hbtheorem}{thm:p17-rat93-ess-restrictions}{proved=annex:p17.rat93-ess-restrictions}{Dominance, support and finite multiplicity}
No weakly dominated strategy is an ESS. Distinct ESSs cannot have the
same support. Hence an \(n\)-strategy matrix game has at most
\(2^n-1\) ESSs and may have none. This does not say that every support
has an ESS or that every finite game has an ESS.
\end{hbtheorem}

These are the chapter's Theorems~1--3, with the rare-mutant and
local-superiority equivalences made explicit. ESS excludes some weakly
dominated Nash equilibria even when they are Pareto efficient.
\cite[Sec.~3, pp.~554--556]{Alexander2021EvolutionaryGames}

\begin{hbexample}{def:p17-rat93-source-static-games}{The source's invasion and welfare contrasts}
Figure~9.3.1a has the row-payoff matrix
\(A=\left(\begin{smallmatrix}2&1\\1&1\end{smallmatrix}\right)\).
Both \((S_1,S_1)\) and \((S_2,S_2)\) are Nash. A rare \(S_1\)
mutant against all \(S_2\) initially ties, but its encounters with
other mutants favor it. The missing strict first-order difference does
not establish invasion resistance.

Figure~9.3.1b is Hawk--Dove:
\[
 A=\begin{pmatrix}(V-C)/2&V\\0&V/2\end{pmatrix},\qquad 0<V<C.
\]
The resource is worth \(V\); the losing Hawk pays injury cost \(C\).
The equal probability of losing a Hawk--Hawk fight yields \((V-C)/2\).
The original animal-conflict terminology was Hawk--Mouse, as Alexander's
note~3 records. The pure Nash profiles are \((H,D)\) and \((D,H)\):
there is no \emph{symmetric} pure Nash profile. The symmetric mixed
profile uses Hawk with probability \(V/C\). It is not strict Nash,
but satisfies the ESS second comparison.

Figure~9.3.2 has
\[
 A=\begin{pmatrix}1&100\\0&100\end{pmatrix}.
\]
Here \(S_1\) weakly dominates \(S_2\), although \((S_2,S_2)\)
is Nash and Pareto efficient with payoffs \((100,100)\). The unique
ESS \(S_1\) has self-play payoff \(1\). Evolutionary stability is
therefore not a welfare criterion.
\end{hbexample}

\begin{hbtheorem}{thm:p17-rat93-static-calculations}{proved=annex:p17.rat93-static-calculations}{Complete calculations for the static contrasts}
For Figure~9.3.1a, \(S_1\) is the only ESS and \(S_2\) is Nash
but not ESS. If the mutant fraction is \(\varepsilon\), its payoff
advantage is \(\varepsilon\).
In Hawk--Dove the complete Nash set is the two asymmetric pure
profiles and the symmetric mixture \(\sigma^*=(V/C,1-V/C)\).
For \(\mu=(q,1-q)\),
\[
 \pi(\sigma^*\mid\mu)-\pi(\mu\mid\mu)
       =\frac C2(q-V/C)^2.
\]
The source's welfare game has exactly the symmetric Nash strategies
\(S_1,S_2\), but only \(S_1\) is ESS. Under replicator dynamics,
its \(S_1\) share \(x\) obeys \(\dot x=x^2(1-x)\). Every
\(0<x(0)<1\) converges to \(1\), while mean payoff
\(100-100x+x^2\) strictly decreases along that trajectory.
\end{hbtheorem}

\section{From revision rules to population dynamics}
\label{sec:p17-rat93-dynamics}

\begin{hbdefinition}{def:p17-rat93-revision}{Revision protocols and their information requirements}
A revision protocol is a locally Lipschitz map
\[
 \rho:\mathbb R^n\times\Delta_n\longrightarrow\mathbb R_+^{n\times n},
\]
where \(\rho_{ij}(\pi,\vec p)\) is the conditional rate of a user
of \(S_i\) switching to \(S_j\). Here the source's repeated glyph
\(\pi\) means a numerical payoff vector, not the earlier bilinear
payoff function. The mean dynamic is the inflow-minus-outflow equation
\[
 \dot p_i=\sum_jp_j\rho_{ji}(W(\vec p),\vec p)
            -p_i\sum_j\rho_{ij}(W(\vec p),\vec p).
\]
Diagonal switch rates cancel. Rates are not arbitrary probabilities:
a discrete revision-opportunity model must choose its opportunity rate
large enough to normalize the conditional switching probabilities.
The mean equation is a deterministic aggregate model, not an exact
sample path of every finite stochastic population.

Write \([a]_+=\max(a,0)\) and
\(\phi(\vec p)=\sum_ip_iW_i(\vec p)\). The three source protocols are
\[
\begin{array}{c|c}
 \text{pairwise proportional imitation}&
 \rho_{ij}=p_j[W_j-W_i]_+\\
 \text{BNN excess-payoff revision}&
 \rho_{ij}=[W_j-\phi]_+\\
 \text{Smith pairwise comparison}&
 \rho_{ij}=[W_j-W_i]_+.
\end{array}
\]
Imitation requires encountering a user of the candidate strategy;
Smith comparison needs its candidate and incumbent payoffs; BNN
comparison also needs the aggregate mean payoff. Comparisons of
expected or well-estimated payoffs are not automatically implemented
by observing one noisy realized match. Alexander's informal learning
interpretation and its payoff-information assumptions must be kept
together.
\end{hbdefinition}

\begin{hbtheorem}{thm:p17-rat93-protocols}{proved=annex:p17.rat93-protocols}{Three derived mean equations and their boundary behavior}
The displayed protocols respectively give
\[
\begin{aligned}
 \dot p_i&=p_i[W_i-\phi] &&\text{(replicator)},\\
 \dot p_i&=[W_i-\phi]_+-p_i\sum_j[W_j-\phi]_+
                                      &&\text{(BNN)},\\
 \dot p_i&=\sum_jp_j[W_i-W_j]_+-p_i\sum_j[W_j-W_i]_+
                                      &&\text{(Smith)}.
\end{aligned}
\]
For locally Lipschitz \(W\) these have unique forward solutions
remaining in \(\Delta_n\) for all finite times. In the replicator
equation, an initially positive share stays positive at every finite
time and an initially zero share stays zero. Extinction can still
occur in the limit. BNN and Smith can introduce a missing strategy.
For either of these two dynamics, rest points are exactly population
Nash states, \(p^{\mathsf T}W(p)\geq q^{\mathsf T}W(p)\) for all
\(q\in\Delta_n\). Replicator rest points need not be Nash.
\end{hbtheorem}

Figure~9.3.5 uses bad Rock--Paper--Scissors,
\[
 A=\begin{pmatrix}0&-2&1\\1&0&-2\\-2&1&0\end{pmatrix}.
\]
At the all-Rock vertex BNN immediately creates Paper, since
\(W_P=1>\phi=0\); when another absent strategy becomes advantageous,
it too can enter. The source's phase portrait depicts a cycling
long-run outcome rather than convergence to an optimal fixed strategy.
An observed orbit in that numerical illustration is not, by itself,
a theorem for every parameter or revision family.
\cite[Sec.~4, pp.~556--559]{Alexander2021EvolutionaryGames}

\begin{hbtheorem}{thm:p17-rat93-replicator-dominance}{proved=annex:p17.rat93-replicator-dominance}{Strict elimination and persistent weak domination}
In a finite linear single-population game, if every strategy has
positive initial share and \(S_i\) is strictly dominated by a pure
or mixed strategy \(\mu\), then replicator dynamics gives
\(p_i(t)\to0\). The conclusion can fail when necessary dominating
strategies are initially absent.

For Figure~9.3.4,
\[
 A=\begin{pmatrix}1&1&1\\1&1&0\\0&0&0\end{pmatrix},
\]
\(S_2\) is weakly dominated by \(S_1\). From every interior state
the replicator solution converges to a mixture with positive shares
of \(S_1\) and \(S_2\), and zero share of \(S_3\).
\end{hbtheorem}

The missing-strategy failure is concrete in the prisoner's dilemma:
if Cooperate is strictly dominated by Defect but the initial population
contains only Cooperate, imitation cannot discover Defect and all-Cooperate
is a rest point. The elimination result concerns a fixed payoff matrix
and a stated initial support, not every plausible learning rule.
\cite[Theorem~4, p.~557]{Alexander2021EvolutionaryGames}

\begin{hbdefinition}{def:p17-rat93-comparison-families}{The actual classes behind the survival statement}
In the primary's family notation, \(\phi_j:\mathbb R\to\mathbb R_+\)
is a scalar response function, a different typed role from the chapter's
population mean \(\phi:\Delta_n\to\mathbb R\).
The broader impartial pairwise class has
\(\rho_{ij}=\phi_j(W_j-W_i)\), with locally Lipschitz
\(\phi_j(d)>0\) exactly when \(d>0\), and \(\phi_j(d)=0\) otherwise.
Impartiality removes dependence on \(i\) from the response function,
not from the incumbent payoff inside its argument. The separable
excess-payoff class uses \(\rho_{ij}=\phi_j(W_j-\phi(p))\), where
the inner \(\phi(p)\) still denotes the population mean.

For the particular one-sixth survival theorem, the primary fixes a
\emph{common} Lipschitz response function \(\phi_j=\phi\), now
with scalar domain \(\mathbb R\), and additionally requires
the positive right derivative \(\phi'_+(0)>0\). The repeated
published glyph is retained with its domain declared, not silently
identified with Alexander's mean-payoff function. BNN and Smith satisfy
these conditions with the response function \(\phi(d)=[d]_+\).
\end{hbdefinition}

\begin{hbtheorem}{thm:p17-rat93-survival}{cited=HofbauerSandholm2011Survival}{Dominated strategies can survive under innovative dynamics}
For either common-function class with the stated sign and derivative
conditions, and for any prescribed positive excluded-volume tolerance,
there is a finite linear game in which a strictly dominated strategy's
share is eventually bounded away from zero and exceeds \(1/6\)
infinitely often, for initial states outside an open set of at most
that tolerance. The game depends on the dynamic and tolerance.
This is not a claim of survival in every game, nor an almost-everywhere
claim for a single game independent of the tolerance.
\end{hbtheorem}

The primary's construction duplicates Scissors in bad RPS and then
reduces the duplicate's payoff by a small \(d>0\):
\[
 \widetilde A_d=
 \begin{pmatrix}
 0&-b&a&a\\ a&0&-b&-b\\ -b&a&0&0\\
 -b-d&a-d&-d&-d
 \end{pmatrix},\qquad b>a>0.
\]
The fourth row is strictly dominated by the third by the uniform gap
\(d\). This is the symmetric normal-form construction underlying
Figure~9.3.6, rather than its inconsistent printed payoff pair.
The proof establishes a repelling equilibrium segment in the identical-twin
game, convergence toward equal twin shares, a cycling attractor, and
persistence of that attractor after a sufficiently small payoff reduction.
Attractor-continuation and the appendix lemmas are the primary's cited
proof boundary. The source's wider summary is not substituted for those
hypotheses. \cite[Sec.~3, Theorem~1 and note~6, pp.~346--351]{HofbauerSandholm2011Survival}

The same primary gives a broader but weaker theorem: Lipschitz
growth rates, positive correlation of motion with payoff away from
rest, Nash stationarity, and innovation of absent optimal strategies
at non-Nash states suffice for a game with persistence from most initial
conditions. That theorem does not retain the one-sixth conclusion
without additional symmetry. These are dynamical disequilibrium effects,
not irrational players being assigned an optimal dominated action at
a Nash equilibrium. \cite[Sec.~4.1, Theorem~2, pp.~351--353]{HofbauerSandholm2011Survival}

\section{Which stability claim is being made?}
\label{sec:p17-rat93-stability}

\begin{hbdefinition}{def:p17-rat93-stability}{Rest, Lyapunov stability, attraction and basins}
A rest point \(\vec p^*\) has \(\dot p_i=0\) for all \(i\).
It is Lyapunov stable, relative to \(\Delta_n\), when for every
\(\varepsilon>0\) some \(\delta>0\) makes
\[
 \lVert\vec p(0)-\vec p^*\rVert<\delta
 \quad\Longrightarrow\quad
 \lVert\vec p(t)-\vec p^*\rVert<\varepsilon\quad(t\geq0).
\]
It is locally attracting when all trajectories from some neighborhood
converge to it. Asymptotic stability requires both Lyapunov stability
and local attraction. Alexander uses \emph{neutrally stable} for a
Lyapunov-stable rest point that is not asymptotically stable; this
dynamic usage is not the separate game-theoretic definition of a
neutrally stable strategy. The basin of attraction is the set of
initial states whose trajectories converge to the point.
No probability distribution on initial conditions is supplied by a
basin's existence or size.
\end{hbdefinition}

\begin{hbtheorem}{thm:p17-rat93-stability-bridges}{proved=annex:p17.rat93-stability-bridges}{The linear replicator bridges, with their directions}
For a finite symmetric two-player game interpreted through random
matching, every symmetric population Nash state is a replicator rest
point. Every Lyapunov-stable replicator rest point is Nash. Every
ESS state is asymptotically stable under replicator dynamics.
These implications neither identify all rest points with Nash nor
identify Lyapunov stability with ESS.
\end{hbtheorem}

\begin{hbexample}{def:p17-rat93-rps}{The source's ESS-free stable cycles}
For ordinary Rock--Paper--Scissors, Figure~9.3.3,
\[
 A=\begin{pmatrix}0&-1&1\\1&0&-1\\-1&1&0\end{pmatrix}.
\]
The unique Nash profile is the uniform mixture for both players.
Its population state \(\vec p^*=(1/3,1/3,1/3)\) is a replicator
rest point and is Lyapunov stable, but is not ESS and is not
asymptotically stable. Interior trajectories off the center cycle
on level curves, as Figure~9.3.8 illustrates.
\end{hbexample}

\begin{hbtheorem}{thm:p17-rat93-rps-calculation}{proved=annex:p17.rat93-rps-calculation}{The invariant explaining neutral rather than asymptotic stability}
For ordinary RPS, \(\phi=0\) and the interior quantity
\(p_1p_2p_3\) is constant along replicator trajectories. It has its
unique maximum \(1/27\) at the uniform state. That invariant gives
Lyapunov stability but excludes convergence to the center from any
other interior state. All three simplex vertices are rest points
but are not Nash. The uniform strategy fails the ESS tie comparison.
\end{hbtheorem}

\begin{hbdefinition}{def:p17-rat93-regular}{Regular ESS in the linear source domain}
A Nash strategy \(\sigma\) is a regular ESS when every unused pure
strategy has strictly smaller payoff against \(\sigma\), and
\[
 x^{\mathsf T}Ax<0
 \quad\text{for every nonzero }x\in\mathbb R^n
 \text{ with }\sum_i x_i=0
 \text{ and }x_i=0\text{ outside }\operatorname{supp}(\sigma).
\]
Support here is contained in, not necessarily properly contained in,
the equilibrium support. This is a negative quadratic form on the
face's tangent space; \(x\) is a signed tangent vector, not a
population distribution. The unused-strategy inequality is often
called quasistrictness.
\end{hbdefinition}

\begin{hbtheorem}{thm:p17-rat93-regular-ess}{proved=annex:p17.rat93-regular-ess}{Regularity is sufficient but not necessary for ESS}
Every regular ESS in this linear domain is an ESS. The pure strategy
\(S_1\) in \(A=\left(\begin{smallmatrix}1&1\\1&0\end{smallmatrix}\right)\)
is an ESS but not regular: \(S_2\) ties against \(S_1\), though
it loses the second ESS comparison.
\end{hbtheorem}

\begin{hbtheorem}{thm:p17-rat93-regular-families}{cited=Sandholm2010LocalStability}{The actual cross-family local stability result}
A regular ESS is asymptotically stable under every locally Lipschitz
sign-preserving impartial pairwise comparison dynamic, every such
separable excess-payoff dynamic, and the best-response dynamic
\(\dot p\in\operatorname{argmax}_{q\in\Delta_n}q^{\mathsf T}W(p)-p\).
For best response the conclusion concerns solutions of a differential
inclusion, not an unspecified unique ordinary differential equation.
The linear case above is included in the primary's continuously
differentiable population-payoff formulation.
\end{hbtheorem}

This is Sandholm's Theorem~1, whose Lyapunov constructions and complete
proof are in its appendix; it is a sufficient \emph{local} condition,
not a global selection theorem. Alexander's attribution to the 2011
survival paper is repaired accordingly. Alexander also reports that,
when a payoff function admits a nonregular ESS, arbitrarily close payoff
functions can be found whose ESSs are all regular. This is an attributed
external perturbation claim, not a locally proved genericity theorem or
a probability statement about real-world games or initial populations.
\cite[Secs.~3--4 and appendix, pp.~33--35, 42--47]{Sandholm2010LocalStability}

\section{Different roles, small basins and unstable global attraction}
\label{sec:p17-rat93-limits}

\begin{hbexample}{def:p17-rat93-two-populations}{Battle of the Sexes with separate role populations}
Figure~9.3.7 uses
\[
 \begin{array}{c|cc}&S_1&S_2\\\hline
 S_1&(1,2)&(0,0)\\ S_2&(0,0)&(2,1)
 \end{array}.
\]
Let \(p_i^R,p_i^C\) be the Row and Column population shares.
Their payoffs are
\[
 W_1^R=p_1^C,\quad W_2^R=2p_2^C,
 \qquad W_1^C=2p_1^R,\quad W_2^C=p_2^R.
\]
The role-specific means
\(\phi^R=p_1^Rp_1^C+2p_2^Rp_2^C\) and
\(\phi^C=2p_1^Rp_1^C+p_2^Rp_2^C\) depend on both populations.
One cannot replace these by a one-population symmetric ESS problem.
The broader extensive-form multipopulation theory cited by the chapter
is a distinct external presentation.
\end{hbexample}

\begin{hbtheorem}{thm:p17-rat93-two-populations}{proved=annex:p17.rat93-two-populations}{All rest points and Nash profiles in the source's two-role example}
Writing \(x=p_1^R\) and \(y=p_1^C\), the coupled replicator equations are
\[
 \dot x=x(1-x)(3y-2),\qquad
 \dot y=y(1-y)(3x-1).
\]
The four corners and \((1/3,2/3)\) are exactly the rest points.
Only \((0,0),(1,1),(1/3,2/3)\) are Nash profiles.
The two matched corners are asymptotically stable; the mixed point is
a saddle with eigenvalues \(2/3,-2/3\). The mismatched corners are
unstable rest points maintained by absence of the improving strategies.
\end{hbtheorem}

\begin{hbexample}{def:p17-rat93-basins}{The source's two final basin contrasts}
Figure~9.3.9 has
\[
 A_\varepsilon=
 \begin{pmatrix}0&1&\varepsilon\\0&1&0\\0&0&2\varepsilon\end{pmatrix},
 \qquad 0<\varepsilon<1/2.
\]
The pure state \(S_3\) is a strict Nash strategy, regular ESS and
asymptotically stable, yet its basin can be made arbitrarily small.
Figure~9.3.10 has
\[
 A=\begin{pmatrix}0&0&2\\1&0&0\\0&0&1\end{pmatrix}.
\]
Every interior replicator trajectory converges to all-\(S_2\), but
all-\(S_2\) is not Lyapunov stable. Near initial states can make a
large excursion before returning. The point is \(S_2\), not the
source prose's undefined \(A\); apparent termination of plotted
trajectories midway through the simplex is explicitly a numerical
solver limitation, not an additional equilibrium.
\end{hbexample}

\begin{hbtheorem}{thm:p17-rat93-small-basin}{proved=annex:p17.rat93-small-basin}{A quantitative bound on the regular ESS basin}
For \(A_\varepsilon\), every state with \(p_2>2\varepsilon\)
remains outside the all-\(S_3\) basin. Consequently that basin is
contained in the strip \(p_2\leq2\varepsilon\), whose relative
two-dimensional simplex area is \(4\varepsilon-4\varepsilon^2\).
It tends to zero as \(\varepsilon\downarrow0\), despite the
regular ESS and local asymptotic stability of \(S_3\).
\end{hbtheorem}

\begin{hbtheorem}{thm:p17-rat93-global-unstable}{proved=annex:p17.rat93-global-unstable}{Full interior attraction without Lyapunov stability}
For Figure~9.3.10, every initially positive distribution converges to
\((0,1,0)\), its unique population Nash state, which is not ESS.
Nevertheless, arbitrarily close feasible states leave
a fixed-size neighborhood of that vertex; interior examples can also
be chosen. Thus a full-interior basin does not imply asymptotic stability.
\end{hbtheorem}

\section{Relations, cognitive limits and consultation}
\label{sec:p17-rat93-consultation}

The classical game interface tests one agent's deviations against fixed
opponents. The epistemic interface
\cref{def:p17-egt-types,def:p17-egt-common-belief}
models subjective conjectures and beliefs about reasoning. Here the
revision protocol describes switching and the population state records
what strategies are actually present. Optimality for beliefs, a Nash
support condition, resistance to mutation, stationarity, local stability,
convergence and welfare are different properties. A population law does
not supply correct beliefs or common belief in rationality; an epistemic
existence theorem does not supply a learning process or its basin.
\cref{thm:p17-strategic-welfare-boundaries} supplies the separate
utility-comparability warning for aggregate welfare.

Evolutionary and cultural interpretations require different mechanisms:
biological inheritance and reproduction are not the same process as
imitation, learning or institutional revision. A payoff must be connected
to the actual transmission mechanism before calling it fitness. The
chapter offers mathematical models and numerical illustrations, not an
experiment estimating how often a named human population uses Smith,
BNN or replicator revision. No study sample, measured frequency,
confidence interval or universal cognitive claim follows from these
equations. Structured matching, heterogeneous information, stochastic
mutation and discrete-time updating require their own models and can
change the results. Normative acceptability, truth and moral obligation
do not follow from a behavior becoming prevalent.
\cite[Secs.~1, 4--5, pp.~553--554, 556--564]{Alexander2021EvolutionaryGames}

\begin{hbmethod}{def:p17-rat93-use}{Match the model to the question before invoking stability}
Fix the agents or role populations, strategy repertoire, payoff meaning,
matching structure and revision information. Distinguish an individual's
mixed strategy from a distribution of pure users. For invasion, check
every first-order best-response tie and its second comparison. For a
trajectory question, specify the actual law and initial support; derive
the mean equation rather than inferring it from the word evolution.
Compute rest points, Nash inequalities and the particular local stability
test separately. For outcome likelihood state the initial-condition
distribution or basin question; for welfare add its distinct evaluative
criterion and utility interpretation.

The Hawk--Dove calculation succeeds at separating mixed ESS from strict
Nash. The welfare game fails any attempted identification of evolutionary
stability with Pareto efficiency or increasing mean payoff. Weak dominance
persists under replicator flow in Figure~9.3.4; strict dominance is
eliminated only with the stated initial support. BNN and Smith can innovate
yet need not eliminate dominated strategies in cycling environments.
RPS separates Lyapunov stability from ESS and attraction. Battle of the
Sexes separates role-specific Nash profiles from boundary rest points.
The final two basin examples defeat both directions of the proposed
identification of local stability with a likely long-run outcome.
\end{hbmethod}

Alexander's figures were produced with the open-source Dynamo package,
as acknowledged in the chapter. Their arrows and shading visualize
specific equations; they are not empirical measurements or independent
proofs. The annex supplies the declared finite derivations. The two
general family results retain explicit cited proof boundaries in the
held primary papers, rather than receiving a fabricated local proof.
